\section{Semantic Vectorization for Input Abstraction}\label{apdx:template}

To enable transfer across schemas, SGIA assigns each abstract bucket $\theta$ a short semantic description $\text{desc}(\theta)$ that is later consumed by the bucket encoder $\mu_b$ (\autoref{sec:vectorization}, \autoref{sec:policy_action}). In our implementation, we use lightweight lexical features (deterministic feature hashing~\cite{weinberger2009featurehashing}) rather than a pretrained language model, which is sufficient to generalize across schemas at low overhead.

We use a schema-agnostic template:

\begin{mdframed}[
    backgroundcolor=gray!10,
    linecolor=white,
    skipabove=0.5em
    ]
\noindent\ttfamily\small
Field=[FieldPath] \\
Type=[SSZType] \\
Aspect=[AspectID] \\
Bucket=[BucketID] \\
Meaning=[NaturalLanguageHint]
\end{mdframed}

\paragraph{What the fields mean}
Each line corresponds to a component of the bucket $\theta$:
\begin{itemize}
    \item \textbf{Field} is the fully-qualified path of the affected schema field (e.g., \icode{BeaconState.JustificationBits}).
    \item \textbf{Type} is the SSZ type of that field (e.g., \icode{Bitvector[4]}).
    \item \textbf{Aspect} identifies \emph{which part of the encoding} is being bucketed (e.g., \icode{Value}, \icode{Length}, \icode{ElementValue}, \icode{Offset}, \icode{Tail}).
    \item \textbf{Bucket} is the bucket identifier within that aspect. For type-driven templates, this ID corresponds to a concrete interval (\autoref{apdx:typp}); for spec-derived boundaries, it is a symbolic label (e.g., \icode{DirtyTail}).
    \item \textbf{Meaning} is a short human-readable gloss of the boundary the bucket represents. It is not required by the concretizer, but provides lexical cues for transfer across schemas and makes buckets auditable.
\end{itemize}

Formally, $\text{desc}(\theta)$ is a deterministic string obtained by concatenating the above components.
The online phase then uses the encoder $\mu_b$ to map this string into a fixed-dimensional vector $h_\theta \in \mathbb{R}^d$.
In our prototype, $\mu_b$ is implemented via deterministic feature hashing~\cite{weinberger2009featurehashing}: for $i\in\{1,\ldots,d\}$,
\[
    h_\theta[i] = {\sf scale}\big(H(\text{desc}(\theta), i)\big) \in [-1,1],
\]
where $H$ is a stable hash function and ${\sf scale}(\cdot)$ maps the hash output into a bounded real value. The full description $\text{desc}(\theta)$ contains all components; the encoder may hash a chosen subset of fields (at minimum Field/Aspect/Bucket) depending on the desired invariances.

In the current artifact, the hashed key includes at least \icode{Field}, \icode{Aspect}, and \icode{Bucket}; \icode{Type} and \icode{Meaning} can be appended without changing the rest of the pipeline.

For example:

\begin{mdframed}[
    backgroundcolor=gray!10,
    linecolor=white,
    skipabove=0.3em,
    skipbelow=0.2em
    ]
\noindent\ttfamily\small
Field=BeaconState.\allowbreak JustificationBits \\
Type=Bitvector[4] \\
Aspect=ElementValue \\
Bucket=B11 \\
Meaning=set non-zero padding bits in the final byte
\end{mdframed}

\begin{mdframed}[
    backgroundcolor=gray!10,
    linecolor=white,
    skipabove=0.2em,
    skipbelow=0.2em
    ]
\noindent\ttfamily\small
Field=Validator.\allowbreak Slashed \\
Type=boolean \\
Aspect=Value \\
Bucket=Dirty \\
Meaning=non-canonical boolean byte in 0x02--0xFF
\end{mdframed}

\begin{mdframed}[
    backgroundcolor=gray!10,
    linecolor=white,
    skipabove=0.2em,
    skipbelow=0.2em
    ]
\noindent\ttfamily\small
Field=SuffixStateDiff.\allowbreak NewValue \\
Type=Optional[Bytes32] \\
Aspect=Tail \\
Bucket=DirtyTail \\
Meaning=append trailing bytes to the None selector
\end{mdframed}

\begin{mdframed}[
    backgroundcolor=gray!10,
    linecolor=white,
    skipabove=0.2em
    ]
\noindent\ttfamily\small
Field=PendingAttestation.\allowbreak AggregationBits \\
Type=Bitlist[2048] \\
Aspect=ElementValue \\
Bucket=B0 \\
Meaning=inject 0x00 to probe the missing-sentinel boundary
\end{mdframed}

\paragraph{Interpreting the examples}
The bucket IDs above correspond to concrete constraints in $\Theta$.
For instance, \icode{B11} refers to the highest byte bucket $[0xF2,0xFF]$ (\autoref{apdx:typp}); when applied to \icode{Bitvector[4]} stored in one byte, sampling from this bucket naturally sets non-zero padding bits.
Similarly, \icode{Dirty} is the boolean non-canonical bucket $[0x02,0xFF]$, \icode{DirtyTail} triggers the union trailing-byte boundary, and \icode{B0} includes \icode{0x00}, which probes the missing-sentinel edge case for bitlists.

In our implementation, these descriptions are embedded using deterministic feature hashing to obtain a fixed-dimensional bucket vector.
This choice keeps the prototype self-contained while still enabling lexical transfer: buckets that share cues such as ``dirty padding'' across schemas map to overlapping hashed features.
